\section{Results}
\label{sec:results}

% This section carries eleven floats in a few pages.  LaTeX's default float
% fractions reserve so much of each page for text that the tables would be
% pushed several pages past the prose that discusses them; the values below let
% a page carry more table and fewer orphaned strips of text.  They are restored
% to the LaTeX defaults at the end of the section.
\renewcommand{\topfraction}{0.92}
\renewcommand{\bottomfraction}{0.75}
\renewcommand{\textfraction}{0.06}
\renewcommand{\floatpagefraction}{0.80}
\setcounter{topnumber}{3}
\setcounter{bottomnumber}{2}
\setcounter{totalnumber}{4}
\setlength{\textfloatsep}{12pt plus 2pt minus 4pt}
\setlength{\floatsep}{10pt plus 2pt minus 3pt}
\setlength{\intextsep}{12pt plus 2pt minus 3pt}

\subsection{Engine and belief validation}
\label{sec:results-verify}

\paragraph{Information-safety audit (E1).}
The verification sweep plays the ordered cross-product of nine policies with the
knowledge audit of \S\ref{sec:engine} active and recorded \numAuditViolations{}
violations in \numAuditChecks{} soundness checks: no derived knowledge state,
capacity or certificate was ever contradicted by the actual deal. Card
conservation, resolution of all nine half-suits and bit-identical replay hold in
every game, and the action-limit rate is \numLimitRate{}. The same run under the
v0.3 dialect records \numAuditViolations{} violations in \numAuditChecksLegacy{}
checks and, again, no action-limit game. E13 measures action-cap incidence over
\numEthirteenDeals{} deals in six seat rotations, seed $464646$, against every
panel member and the \vfast{} mirror, and records \numLimitRate{} on every row.

\paragraph{Belief cross-checks (E2).}
Over \numSelftestGames{} games and \numSelftestChecks{} comparisons the exact
reference engine agrees with the card-level capacity dynamic program to
$\numBeliefVsCardDP$ with no certificate live --- machine precision, as it must,
the two computing one quantity two ways --- and with exact rejection sampling to
$\numBeliefVsSampling$ with certificates live, Monte-Carlo error at these sample
sizes. The Fast approximation has mean absolute marginal error $\numSinkhornMean$
and maximum $\numSinkhornMax$: accurate on average, with a heavy tail that E2
does not localise to any class of position.

\paragraph{Brute-force allocation oracle (E15).}
E2 compares two structured algorithms against each other and does not establish
that either enumerates the correct collection of assignments. E15 tests that
directly against an enumerator using no dynamic programming, no factorisation and
no sampling; the procedure, its per-state cut-off of \numOracleMaxDeals{}
consistent deals and the sampler comparison are in \S\ref{app:inference}.


Every deterministic comparison (Table~\ref{tab:oracle}, \S\ref{app:oracle}) differs by exactly zero,
which is exact rather than approximate agreement because both sides are integer
counts in double precision; of \numOracleBestChecks{} checks on the chosen
allocation \numOracleBestBad{} were inconsistent and none failed to be the argmax
of $\alpha$, and sampler frequencies fall within \numOracleSampleDiff{} of the
exact marginals over \numOracleDraws{} draws. Coverage is partial by construction: \numOracleStatesSkipped{} states,
more than were enumerated, were skipped as too large, and skipping is correlated
with what makes a state hard, a large space of unresolved cards. E15 validates
the reference engine on small reachable states, not the whole state space.

\paragraph{Declaration pre-gate audit (E16).}
Two cheap pre-gates screen declaration candidates before the full posterior
evaluation runs (\S\ref{sec:declare}), and neither is proved to be an upper bound
on the quantity it screens. E16 re-runs the full evaluation and the same stopping
rule with both gates disabled, on every live half-suit at every declaration
opportunity of the primary bank (\S\ref{app:ablations}).


Of \numGateSeen{} (opportunity, live half-suit) pairs over \numGateOpps{}
opportunities, \numGateRejected{} (\numGateRejectedPct\%) were gated out, and
\numGateFalseNeg{} of those --- \numGateFalseNegPct\% --- would have been
declared by the full rule; the gated path makes \numGateDeclGated{} declarations
against \numGateDeclUngated{} ungated, choosing differently at
\numGateActionsChanged{} opportunities, \numGateActionsPct\% of the total and
\numGateDeclLostPct\% of the ungated rule's declarations
(Table~\ref{tab:gateaudit}, \S\ref{app:abl-gateaudit}). The pre-gates are
therefore not lossless: they are a pruning heuristic with a measured
false-negative rate of roughly one in ten thousand declaration opportunities,
specific to this policy, this corpus and the compiled thresholds. The audit
bounds the size of that effect rather than eliminating it, and does not license
describing the gates as safe or as a proved upper bound.

\paragraph{Port validation (E8).}
Every comparison in this paper is made against a port of \vthree{} into the new
engine. Under the v0.3 dialect that port reproduces six of the seven published
figures within their intervals (Table~\ref{tab:port}, \S\ref{app:port}). The row for the
misdirection artist misses by half a point, near the top of a scale on which both
figures exceed $98\%$.


\subsection{Primary result: \vfast{} against \vthree{}}
\label{sec:results-h2h}

On the primary evaluation bank of \numHtHDeals{} held-out deals played in six
seat rotations for \numHtHGames{} games, \vfast{} defeated \vthree{} in
\numVsVthreeWins{} games --- a \numVsVthree\% win rate with a deal-clustered
$95\%$ interval of \numVsVthreeCI\% --- taking \numVsVthreeSets{} of the nine
half-suits on average.


Table~\ref{tab:elo} converts the round-robin --- reproduced cell by cell as
Table~\ref{tab:matrix} in \S\ref{app:pop-matrix} --- into
Bradley--Terry ratings: \vfast{} is rated \numElo{} on a scale where \vthree{} is
zero, against \numEloDetective{} for the posterior detective and \numEloLockout{}
for the turn-starvation lockout, so those two sit within a few Elo of \vthree{}.
Such ratings are not invariant to redundancy in the population and cannot
represent cyclic dominance \cite{balduzzi-reeval}, so the table summarises this
panel rather than giving a transferable strength scale.

\subsection{Calibration}
\label{sec:results-calib}

\begin{figure}[!htbp]
\centering
\resizebox{\linewidth}{!}{% Generated by engine/build_tables.py -- do not edit by hand.
\begin{tikzpicture}
\begin{scope}
\draw[fishgray!35] (0,0) grid[step=0.5000] (5.000,5.000);
\draw[fishgray!70,thick] (0,0) rectangle (5.000,5.000);
\draw[fishgray,dashed] (0,0) -- (5.000,5.000);
\node[below,font=\scriptsize,text=fishgray] at (0.000,0) {0.0};
\node[left,font=\scriptsize,text=fishgray] at (0,0.000) {0.0};
\node[below,font=\scriptsize,text=fishgray] at (2.500,0) {0.5};
\node[left,font=\scriptsize,text=fishgray] at (0,2.500) {0.5};
\node[below,font=\scriptsize,text=fishgray] at (5.000,0) {1.0};
\node[left,font=\scriptsize,text=fishgray] at (0,5.000) {1.0};
\node[below,font=\scriptsize,text=fishgray] at (2.500,-0.42) {forecast};
\node[rotate=90,above,font=\scriptsize,text=fishgray] at (-0.45,2.500) {observed};
\node[above,font=\small\bfseries,text=fishgreen] at (2.500,5.100) {v0.4-Fast: $\Pr[\text{ask succeeds}]$};
\fill[fishgreen,opacity=0.80] (0.1455,0.0630) circle (0.1407);
\fill[fishgreen,opacity=0.80] (0.7785,0.7035) circle (0.1334);
\fill[fishgreen,opacity=0.80] (1.2245,1.4380) circle (0.1434);
\fill[fishgreen,opacity=0.80] (1.7415,2.0500) circle (0.1415);
\fill[fishgreen,opacity=0.80] (2.2480,2.5285) circle (0.1353);
\fill[fishgreen,opacity=0.80] (2.7010,2.8955) circle (0.1332);
\fill[fishgreen,opacity=0.80] (3.2230,3.4590) circle (0.1175);
\fill[fishgreen,opacity=0.80] (3.7610,3.5185) circle (0.1164);
\fill[fishgreen,opacity=0.80] (4.2480,4.1025) circle (0.1090);
\fill[fishgreen,opacity=0.80] (4.9900,4.9815) circle (0.1684);
\draw[fishgreen,thick,opacity=0.55] (0.1455,0.0630) -- (0.7785,0.7035) -- (1.2245,1.4380) -- (1.7415,2.0500) -- (2.2480,2.5285) -- (2.7010,2.8955) -- (3.2230,3.4590) -- (3.7610,3.5185) -- (4.2480,4.1025) -- (4.9900,4.9815);
\node[below right,font=\scriptsize,text=fishgray,align=left] at (0.12,4.880) {$n=56{,}204$\\ECE $0.0288$};
\end{scope}
\begin{scope}[xshift=6.600cm]
\begin{scope}
\draw[fishgray!35] (0,0) grid[step=0.5000] (5.000,5.000);
\draw[fishgray!70,thick] (0,0) rectangle (5.000,5.000);
\draw[fishgray,dashed] (0,0) -- (5.000,5.000);
\node[below,font=\scriptsize,text=fishgray] at (0.000,0) {0.0};
\node[left,font=\scriptsize,text=fishgray] at (0,0.000) {0.0};
\node[below,font=\scriptsize,text=fishgray] at (2.500,0) {0.5};
\node[left,font=\scriptsize,text=fishgray] at (0,2.500) {0.5};
\node[below,font=\scriptsize,text=fishgray] at (5.000,0) {1.0};
\node[left,font=\scriptsize,text=fishgray] at (0,5.000) {1.0};
\node[below,font=\scriptsize,text=fishgray] at (2.500,-0.42) {forecast};
\node[rotate=90,above,font=\scriptsize,text=fishgray] at (-0.45,2.500) {observed};
\node[above,font=\small\bfseries,text=fishgreen] at (2.500,5.100) {v0.4-Fast: $\Pr[\text{allocation correct}]$};
\fill[fishgreen,opacity=0.80] (0.0000,0.0000) circle (0.0729);
\fill[fishgreen,opacity=0.80] (2.6695,2.6135) circle (0.0788);
\fill[fishgreen,opacity=0.80] (3.3120,4.1215) circle (0.0775);
\fill[fishgreen,opacity=0.80] (3.7170,4.8315) circle (0.0928);
\fill[fishgreen,opacity=0.80] (4.2065,4.7060) circle (0.0863);
\fill[fishgreen,opacity=0.80] (4.9945,4.9930) circle (0.1347);
\draw[fishgreen,thick,opacity=0.55] (0.0000,0.0000) -- (2.6695,2.6135) -- (3.3120,4.1215) -- (3.7170,4.8315) -- (4.2065,4.7060) -- (4.9945,4.9930);
\node[below right,font=\scriptsize,text=fishgray,align=left] at (0.12,4.880) {$n=5{,}816$\\ECE $0.0222$};
\end{scope}
\end{scope}
\begin{scope}[xshift=13.200cm]
\begin{scope}
\draw[fishgray!35] (0,0) grid[step=0.5000] (5.000,5.000);
\draw[fishgray!70,thick] (0,0) rectangle (5.000,5.000);
\draw[fishgray,dashed] (0,0) -- (5.000,5.000);
\node[below,font=\scriptsize,text=fishgray] at (0.000,0) {0.0};
\node[left,font=\scriptsize,text=fishgray] at (0,0.000) {0.0};
\node[below,font=\scriptsize,text=fishgray] at (2.500,0) {0.5};
\node[left,font=\scriptsize,text=fishgray] at (0,2.500) {0.5};
\node[below,font=\scriptsize,text=fishgray] at (5.000,0) {1.0};
\node[left,font=\scriptsize,text=fishgray] at (0,5.000) {1.0};
\node[below,font=\scriptsize,text=fishgray] at (2.500,-0.42) {forecast};
\node[rotate=90,above,font=\scriptsize,text=fishgray] at (-0.45,2.500) {observed};
\node[above,font=\small\bfseries,text=fishgreen] at (2.500,5.100) {v0.3: $\Pr[\text{allocation correct}]$};
\fill[fishgreen,opacity=0.80] (3.9475,0.9675) circle (0.0925);
\fill[fishgreen,opacity=0.80] (4.2295,1.6705) circle (0.0938);
\fill[fishgreen,opacity=0.80] (4.8610,4.6730) circle (0.1304);
\draw[fishgreen,thick,opacity=0.55] (3.9475,0.9675) -- (4.2295,1.6705) -- (4.8610,4.6730);
\node[below right,font=\scriptsize,text=fishgray,align=left] at (0.12,4.880) {$n=4{,}984$\\ECE $0.1257$};
\end{scope}
\end{scope}
\end{tikzpicture}
}
\caption{Reliability diagrams (E6) for the two forecasts \vfast{} exposes and for
\vthree{}'s declaration forecast. Population: the calibration bank of
\numCalibDeals{} deals, seed $717171$, \vfast{} against \vthree{}, weights
frozen; the unit is one forecast. Marker size increases with bin count on a
compressed scale, so it orders the bins rather than measuring them.
Points above the diagonal are underconfident, below it overconfident. No
clustered intervals are computed.}
\label{fig:reliability}
\end{figure}

\begin{table}[!htbp]
\centering
\caption{Aggregate forecast reliability (E6). Population: the calibration bank of
\numCalibDeals{} deals, seed $717171$, \vfast{} against \vthree{}, both frozen;
the unit is one forecast and $n$ counts forecasts of that type. \vthree{} exposes
no ask forecast, so only the declaration rows are comparable. All entries are
point estimates: no clustered uncertainty for Brier score, log loss or ECE.
Voluntary and forced declarations are pooled here, unlike in
Table~\ref{tab:h2h}.}
\label{tab:calib}
\small
\begin{tabular}{llrrrrr}
\toprule
Policy & Forecast & $n$ & Brier & Log loss & ECE & Mean pred / obs \\
\midrule
FishBot v0.4 & $\Pr[\text{ask succeeds}]$ & 56{,}204 & 0.1297 & 0.3790 & 0.0288 & 0.5338 / 0.5504 \\
FishBot v0.4 & $\Pr[\text{allocation correct}]$ & 5{,}816 & 0.0155 & 0.0592 & 0.0222 & 0.9575 / 0.9789 \\
\midrule
FishBot v0.3 & $\Pr[\text{allocation correct}]$ & 4{,}984 & 0.1339 & 0.3990 & 0.1257 & 0.9459 / 0.8202 \\
\bottomrule
\end{tabular}
\end{table}

Table~\ref{tab:calib} gives the aggregate scores, Figure~\ref{fig:reliability}
the reliability diagrams and \S\ref{app:calibration} the decile breakdown. Every
calibration score in this paper --- here, in the appendix, and wherever one is
quoted elsewhere --- is a point estimate with no clustered uncertainty, so small
differences between the estimators are not resolved by these numbers.

\vfast{}'s declaration aggregate is dominated by one bin: \numDeclTopBinShare\% of
the \numDeclN{} forecasts fall in the top decile, close to $1$ and almost exactly
right, which is most of what an expected calibration error of \numDeclECE{}
against realised accuracy \numDeclObs\% measures. The decision-relevant bins are
the middle ones, and in $[0.6, 0.9)$ the forecast is \emph{under}confident ---
realised accuracy exceeds predicted probability by $10.0$ to $22.3$ points --- on counts
of 74, 415 and 221, individually noisy and carrying little aggregate weight.
Whatever the stopping rule of \S\ref{sec:stopping} does there, the aggregate
score does not characterise it. \vthree{} is the contrasting case on the same
bank, expected calibration error \numVthreeDeclECE{} running in the overconfident
direction, mean predicted \numVthreeDeclPred\% against realised
\numVthreeDeclObs\%. That is a descriptive comparison of two estimators on the
same games, not a causal explanation of either policy's behaviour, and it does
not establish why the declaration-accuracy difference of \S\ref{sec:results-h2h}
exists.

\subsection{Which mechanisms matter: frozen-policy ablations}
\label{sec:ablations}

Each ablated variant plays the same \numAblDeals{} deals in two seat rotations
against the same four-policy panel with the fitted weights left \emph{unchanged},
so each row measures how the fitted policy behaves when one component is removed
from it, not the standalone causal value of that component in a policy refitted
around its absence: a cheap removal may still be doing work a refitted policy
would find another way to do, and an expensive removal may be expensive only
because the surrounding weights were fitted to rely on it. No matched-budget
refit was run (\S\ref{sec:limitations}). That qualification applies to every row
below and to \S\ref{app:ablations} and is not repeated per row; all ablation
intervals are paired percentile bootstrap over deals, \numBootstrapDraws{}
resamples.

\begin{table}[!htbp]
\centering
\caption{Frozen-policy ablations (E5). Population: \numAblDeals{} deals in two
seat rotations against the four-policy panel \{\vthree{}, turn-starvation
lockout, posterior detective, \vtwo{}\}, \numAblGames{} games per opponent, seed
$606060$, full rules dialect. Fitted weights are frozen at the shipping values in
every row; only the named component changes. ``Full $-$ ablated'' is the
reference win rate \numAblRef\% minus the row's, so positive means removal hurt.
Intervals: paired percentile bootstrap over deals. The three belief rows at
$21$--$29\%$ substitute weaker approximate beliefs, not the exact reference
engine; only the \vblock{} row uses that. The two policy-prior rows carry
identical figures by construction rather than by coincidence, because the Fast
belief is already the default (\S\ref{app:ablations}).}
\label{tab:ablations}
\footnotesize
\begin{tabular}{@{}l>{\raggedright\arraybackslash}p{0.40\linewidth}rrr@{}}
\toprule
Group & Change (policy weights frozen) & Win rate & Full $-$ ablated & 95\% paired CI \\
\midrule
& \emph{\vfast{} (reference configuration)} & 77.50\% & --- & --- \\
\midrule
Belief & Drop capacities (C4) and the certificates (C5) and the prior; uniform over the surviving support & 21.48\% & +56.03 & +54.27--+57.75 \\
 & Drop the certificates (C5) and the prior; Sinkhorn over C1--C4 & 24.73\% & +52.78 & +50.98--+54.55 \\
 & Drop the certificates (C5) and the prior; exact capacity dynamic program over C1--C4 & 28.60\% & +48.90 & +47.00--+50.78 \\
 & \vblock{} belief substituted into the Fast-fitted policy (exact C1--C5, no prior) & 71.30\% & +6.20 & +4.35--+8.05 \\
 & Fast belief with the policy prior off (matches the reference engine's prior-free belief) & 73.62\% & +3.88 & +2.08--+5.73 \\
 & Remove the fitted policy-aware prior & 73.62\% & +3.88 & +2.08--+5.73 \\
\addlinespace
Ask & Zero the lock-completion weight & 74.30\% & +3.20 & +1.55--+4.85 \\
 & Zero the hit-probability weight & 75.95\% & +1.55 & -0.25--+3.35 \\
 & Remove the two-ply refinement & 76.28\% & +1.23 & -0.53--+3.05 \\
 & Zero both information-leak weights & 77.38\% & +0.12 & -1.38--+1.65 \\
 & Remove the one-ply value lookahead & 77.83\% & -0.33 & -1.88--+1.23 \\
 & Zero the reply-threat weight & 77.85\% & -0.35 & -1.98--+1.27 \\
 & Zero the runway weight & 77.85\% & -0.35 & -1.65--+0.97 \\
\addlinespace
Declaration & Declaration threshold 0.99 & 77.30\% & +0.20 & -0.20--+0.60 \\
 & Name the allocation by conditioned greedy MAP & 77.33\% & +0.18 & -0.12--+0.47 \\
 & Fixed threshold in place of the value-based stopping rule & 77.38\% & +0.12 & -1.23--+1.47 \\
 & Declaration threshold 0.80 & 77.48\% & +0.03 & -0.18--+0.22 \\
 & Cash locked half-suits immediately (inert while the stopping rule is active) & 77.50\% & +0.00 & +0.00--+0.00 \\
\bottomrule
\end{tabular}
\end{table}

\paragraph{Belief.}
The three rows that degrade the belief itself are the largest in
Table~\ref{tab:ablations} by an order of magnitude over anything else, and none
of them isolates a single constraint. The fitted prior of \eqref{eq:prior} is
applied only along the Fast path (\filepath{engine/src/belief.hpp}), so every one
of these rows removes the prior together with whatever else it removes. Keeping
the Sinkhorn solver and dropping the certificates and the prior costs
\numAblSinkhorn{} points (CI \numAblSinkhornCI{}), leaving \numAblSinkhornAbs\%;
making the same removal with the exact capacity dynamic program in place of
Sinkhorn costs \numAblCFive{} points (CI \numAblCFiveCI{}), leaving
\numAblCFiveAbs\%; dropping capacity conditioning as well costs \numAblCFour{}
points (CI \numAblCFourCI{}), leaving \numAblCFourAbs\% and the largest decrease
in the table. The prior alone is worth \numAblPrior{} points, so the certificates
account for most but not all of each difference, and no row in this experiment
separates the two. None of the three uses the exact reference engine --- all
substitute a degraded approximate belief --- and only the \vblock{} row
substitutes the engine of \S\ref{sec:belief}.

\paragraph{Prior and ask features.}
Removing the fitted policy-aware prior costs \numAblPrior{} points (CI
\numAblPriorCI{}), leaving \numAblPriorAbs\%; the fitted policy depends on it. Of
the ask-side rows, which probe single coordinates of the \numAskFeats-weight ask
score and the lookahead terms, only the lock-completion weight has an interval
excluding zero: \numAblLockWeight{} points (CI \numAblLockWeightCI{}), leaving
\numAblLockWeightAbs\%. Every other ask-side interval contains zero --- hit
probability \numAblHitWeight{} (CI \numAblHitWeightCI{}), two-ply
\numAblTwoPly{} (CI \numAblTwoPlyCI{}), information leak \numAblLeak{} (CI
\numAblLeakCI{}) --- and three are negative point estimates: one-ply value
\numAblValue{} (CI \numAblValueCI{}), reply threat \numAblThreat{} (CI
\numAblThreatCI{}), runway \numAblRunway{} (CI \numAblRunwayCI{}). Correlated
features produce exactly this pattern, one coordinate's removal being absorbed by
the others, and a paired design controls the deals but not the collinearity; we
keep all three, because dropping a component on a non-significant ablation is
selection on noise.

\paragraph{Declaration.}
Every declaration-side row is unresolved from zero: the value-based stopping rule
against a fixed threshold at \numAblStop{} points (CI \numAblStopCI{}), the
threshold settings $0.80$ and $0.99$ at \numAblDeclEighty{} (CI
\numAblDeclEightyCI{}) and \numAblDeclNinetyNine{} (CI \numAblDeclNinetyNineCI{}).
The ``cash locked half-suits immediately'' row is zero by construction rather than
by measurement: the value-based rule decides first, so the cashing rule never
fires. With the belief rows and the declaration-accuracy difference of
\S\ref{sec:results-h2h}, these support only the claim that the \emph{joint}
allocation probability matters, which is what the certificates and the capacity
conditioning sharpen. They do not establish that the stopping formulation of
\S\ref{sec:stopping} outperforms a well-calibrated fixed threshold; this
experiment lacks the power to resolve that.

\paragraph{Belief substitution.}
Substituting the \vblock{} belief --- the exact reference engine --- into the
Fast-fitted policy costs \numAblBlock{} points (CI \numAblBlockCI{}), leaving
\numAblBlockAbs\%. The matched comparator --- the Fast belief with the policy
prior switched off, so that it is prior-free as the reference engine is --- costs
\numAblPrior{} points on its own. These are two separate frozen-policy measurements and do not
decompose the \numAblBlock{}-point difference additively; the remainder is in the
same direction and no interval is reported for it. Repeating the substitution
under a deliberately unfitted posterior-greedy policy (E12, \numETwelveDeals{}
deals in six seat rotations, \numETwelveGames{} games per arm), which removes the
confound of weights fitted around Fast beliefs, gives \numETwelveFast\% (CI
\numETwelveFastCI\%) against \numETwelveBlock\% (CI \numETwelveBlockCI\%),
deal-clustered percentile bootstrap: same direction, overlapping intervals, so
E12 does not resolve the question either. The fitted policy is therefore not
invariant to the belief representation. This does not establish that an
exact-belief policy would be weaker after matched retraining, which would require
separate matched-budget fits that we did not run; as an analogy only,
\cite{rebstock-inference} report a comparable observation in Skat, where more
accurate inference did not translate into better play under a policy tuned around
the cheaper estimate.

\subsection{Robustness: rules dialects and arbitration}
\label{sec:robustness}

The margin over \vthree{} depends on the dialect (E7, \numEsevenDeals{} deals in
six seat rotations, \numEsevenGames{} games per row, deal-clustered percentile
bootstrap, tables in \S\ref{app:dialects}): it falls to \numEsevenVthreeLegacy\%
(CI \numEsevenVthreeLegacyCI\%) under the v0.3 dialect, which permits declarations
only on turn, to \numEsevenVthreeEight\% (CI \numEsevenVthreeEightCI\%) under the
$48$-card, eight-half-suit form, and to \numEsevenVthreeNoOOT\% (CI
\numEsevenVthreeNoOOTCI\%) with out-of-turn declarations disabled. Comparator rows
do not move uniformly downwards --- against the lockout, disabling out-of-turn
declarations \emph{raises} the win rate to \numEsevenLockoutNoOOT\% --- and since
E7 and E3 differ in bank as well as dialect, these rows bound the size of the
dialect dependence rather than isolating the rule's contribution.

Arbitration of simultaneous voluntary declarations is a modelling choice rather
than a rule of the game (\S\ref{sec:rules}). Replacing the default lowest-seat
order by highest seat or by a scan from the turn-holder moves every win rate by
at most \numArbMaxMove{} points, with overlapping deal-clustered percentile
bootstrap intervals throughout (E17, \numEseventeenDeals{} deals in six seat
rotations, tables in \S\ref{app:dialects}); against \vthree{} the alternatives
give \numArbHighVthree\% and \numArbTurnVthree\% for a default of \numVsVthree\%,
so both put the worst panel result marginally below the \numWorstCase\% of
\S\ref{sec:results-h2h}, which is why that figure is scoped to the default order.
What the choice does move is the out-of-turn declaration rate, from
\numArbOOTSeatLo--\numArbOOTSeatHi{} per game under either seat order to
\numArbOOTTurnLo--\numArbOOTTurnHi{} under the turn-holder scan, declarations per
game unchanged. Both checks are within-dialect and within-panel.

\subsection{Local response probe}
\label{sec:lbr}

\begin{table}[!htbp]
\centering
\caption{Local response probe (E10). Each row fits a fresh policy in the same
policy class, with the same optimiser and budget, whose only objective is to beat
the frozen policy named; fitting used seed $515253$, evaluation a separate bank
of \numLBRDeals{} deals in six seat rotations, \numLBRGames{} games per row, seed
$6543210$, target weights frozen. The win rate is what the response achieved, so
lower is better for the frozen policy. Intervals: deal-clustered percentile
bootstrap, \numBootstrapDraws{} resamples. Action-limit rate \numLBRLimit{} in
all three probes.}
\label{tab:lbr}
\small
\begin{tabular}{lrrr}
\toprule
Frozen target & Response win rate & 95\% CI & Mean half-suits conceded \\
\midrule
\vfast{} & 51.19\% & 49.67--52.72 & 4.550 \\
\vthree{} & 77.31\% & 75.97--78.67 & 5.547 \\
Posterior detective & 76.47\% & 75.11--77.81 & 5.471 \\
\bottomrule
\end{tabular}
\end{table}

A response fitted specifically against frozen \vfast{} achieved \numLBRFour\% (CI
\numLBRFourCI\%, Table~\ref{tab:lbr}); the interval includes $50\%$, so this
search did not resolve an advantage in either direction. The identical procedure
achieved \numLBRThree\% (CI \numLBRThreeCI\%) against \vthree{} and
\numLBRDetective\% against the posterior detective, so it is sensitive enough to
find large exploits when they exist; the action-limit rate is \numLBRLimit{} in
all three probes, the \vfast{} mirror included. As \S\ref{sec:protocol-metrics}
sets out, an achieved response score is a \emph{lower} bound on the
exploitability of the frozen target within the searched class \cite{lisy-lbr}: a
restricted-policy-class diagnostic, not evidence of an equilibrium.

\subsection{Throughput}
\label{sec:throughput}

\begin{table}[!htbp]
\centering
\caption{Whole-game throughput (E9). Population: self-play games under the full
rules dialect with policy weights frozen; the unit is one complete game and the
last column gives how many were timed. The \vfast{} and \vblock{} rows differ
only in the belief path. No intervals; these are wall-clock rates for one timing
run each.}
\label{tab:throughput}
\small
\begin{tabular}{lrr}
\toprule
Configuration & Games/s & Games timed \\
\midrule
\vfast{} (deployed) & 162.8 & 400 \\
\vblock{} (exact reference belief) & 11.9 & 120 \\
\vthree{} & 6{,}954 & 800 \\
\bottomrule
\end{tabular}
\end{table}

\vfast{} runs at \numThroughputFast{} games/s against \numThroughputBlock{} for
\vblock{}, a factor of about \numThroughputRatio$\times$; \vthree{}, whose
posterior carries neither capacity conditioning nor ask-legality certificates,
runs at \numThroughputThree{} games/s. That Fast-to-Block factor is why the exact
reference engine validates probabilities and serves as an ablation arm rather
than running in the deployed inner loop: at \numThroughputBlock{} games/s the
primary bank of \numHtHGames{} games per row would dominate the experimental
budget. Every evaluation run reported in \S\ref{sec:results} completes in
minutes at these rates; the fitting rounds of \S\ref{sec:fitting} and the
exhaustive enumeration of E15 do not, and their costs are not reported here.

\renewcommand{\topfraction}{0.7}
\renewcommand{\bottomfraction}{0.3}
\renewcommand{\textfraction}{0.2}
\renewcommand{\floatpagefraction}{0.5}
\setcounter{topnumber}{2}
\setcounter{bottomnumber}{1}
\setcounter{totalnumber}{3}
\setlength{\textfloatsep}{20pt plus 2pt minus 4pt}
\setlength{\floatsep}{12pt plus 2pt minus 2pt}
\setlength{\intextsep}{12pt plus 2pt minus 2pt}
